% Blocchi della catena nei PDF: la tesi in apertura, le schede delle fonti,
% la controtesi. Nessun pacchetto in piu': filetti e rientri, perche' devono
% reggere anche su una stampante casalinga.
\providecommand{\ccBadge}[1]{}
% Come sul sito: testo maiuscolo in Archivo stretto, non una pillola.
\renewcommand{\ccBadge}[1]{{\ccArchivo\scriptsize\strut\textcolor{ccblu}{\MakeUppercase{#1}}}}

% La tesi: l'affermazione di cui la catena risponde, in apertura.
\newenvironment{tesi}{%
  \par\addvspace{4pt}%
  \noindent\textcolor{ccblu}{\rule{\linewidth}{0.7pt}}\par\addvspace{8pt}%
  \begingroup\setlength{\parindent}{0pt}%
  {\ccArchivo\scriptsize\textcolor{ccblu}{TESI}\par}\addvspace{4pt}%
}{%
  \par\endgroup\addvspace{8pt}%
  \noindent\textcolor{ccblu}{\rule{\linewidth}{0.7pt}}\par\addvspace{15pt}%
}

% Una fonte incastonata nel testo: rientrata, in corpo minore, fra due filetti.
\newenvironment{scheda}{%
  \par\addvspace{12pt}\begingroup
  \leftskip=9mm\relax\setlength{\parindent}{0pt}\footnotesize\color{ccgrigio}%
  \noindent\textcolor{ccblu!35}{\rule{\dimexpr\linewidth-9mm\relax}{0.4pt}}\par\addvspace{6pt}%
}{%
  \par\addvspace{6pt}%
  \noindent\textcolor{ccblu!35}{\rule{\dimexpr\linewidth-9mm\relax}{0.4pt}}\par
  \endgroup\addvspace{12pt}%
}

% La controtesi: rientrata come la scheda ma senza filetti, per distinguerla.
\newenvironment{controtesi}{%
  \par\addvspace{12pt}\begingroup
  \leftskip=9mm\relax\setlength{\parindent}{0pt}\footnotesize\color{ccgrigio}%
}{%
  \par\endgroup\addvspace{12pt}%
}
